\thispagestyle{empty}\vspace*{\fill}
\begin{minipage}{128mm}\fontsize{9}{13}\selectfont\color{otMuted}
{\otSemi\color{otInk}Mathématiques · Quatrième · France}\par\medskip
Conception et mise en page : collection OptiTAB.\par
270 exercices de chapitre, 18 défis et deux bilans annuels originaux, corrigés.\par
Charte graphique de la collection OptiTAB, du collège au lycée.\par\bigskip
Édition du 22 octobre 2026 · Année scolaire 2026-2027.\par\bigskip
Cette édition suit les attendus de Quatrième applicables en 2026-2027.
Le nouveau programme publié en mars 2026 ne s’appliquera en Quatrième
qu’à la rentrée 2027-2028. Les repères de Cinquième nécessaires sont réactivés
au fil des chapitres.\par\bigskip
L’ordre, les exercices et les barèmes sont des choix éditoriaux. Les défis sont
facultatifs. Les deux bilans sont des évaluations d’entraînement originales,
sans statut d’examen officiel.\par\bigskip
Les données des situations sont construites pour apprendre ; ce ne sont pas
des relevés économiques ou scientifiques. Sauf mention contraire, les
schémas de principe ne sont pas à mesurer. Les coordonnées et graduations
explicites font foi dans les graphiques.\par\bigskip
Conserver les résultats exacts, sauf demande d’arrondi. Les réponses ajoutées
en vert dans les figures corrigées ne sont pas des données de l’énoncé.\par\bigskip
© OptiTAB, 2026.\par\otSite
\end{minipage}\clearpage
